\paragraph{Proposition B.6.17 (Proved: the actual quotient Pontryagin class).}
Use the fixed-circle replacement of B.6.16, the positive parameter
loop, and the mapping-torus convention
$(z,1)\sim(\Phi z,0)$ of B.6.1.  Let $M_s$ be its
seven-dimensional mapping torus and let
\[
 \pi:M_s\longrightarrow Q=M_s/S^1_\delta
\]
be the principal circle bundle, with Euler class $x$.
Orient $Q$ by putting the positive $\delta$ direction first
in the orientation of $M_s$.  The bundle over the parameter
circle has fibre $S^2\times S^3$.  The class $x$ is primitive,
and there is a unique $y\in H^4(Q;\mathbb Z)$ with
$\langle xy,[Q]\rangle=1$.  If $Z_1,Z_2\subset Q$ are the
quotients of the two swept critical elliptic curves, oriented
by their complex normal bundles, then
\[
 p_1(Q)=-2\operatorname{PD}_Q(Z_1+Z_2),\qquad
 \langle x,[Z_1]\rangle=2,\quad
 \langle x,[Z_2]\rangle=0.
 \tag{B.Q1}
\]
Consequently
\[
                       p_1(Q)=-4y.
 \tag{B.Q2}
\]
This computes the integer $\ell$ in B.6.15, with these
orientations.  It does not compute $k$ in $x^2=ky$ or the
relative Pontryagin number of the parametrized surgery trace.

\begin{proof}
The quotient of each six-dimensional fibre is
$B_5=S^2\times S^3$, by the free-circle description in
D.2.14 applied to the nullsphere surgery.  Its circle Euler
class is primitive.  One can see its generator geometrically:
the vanishing $S^3$ of the local conifold smoothing is
invariant under the free circle, and its quotient is the new
$S^2$ in $B_5$.  The restricted circle bundle is a Hopf bundle,
so its Euler number has absolute value one.  In particular
this sphere generates $H_2(B_5;\mathbb Z)$.
The return preserves the circle orientation and hence its
Euler class.  The integral Wang sequence, with
$H^1(B_5)=0$, gives
\[
 H^2(Q;\mathbb Z)\xrightarrow{\;\cong\;}H^2(B_5;\mathbb Z).
 \tag{B.Q3}
\]
It also gives $H^4(Q;\mathbb Z)=\mathbb Z$ and no torsion
in these degrees.  Integral duality proves the assertion
about $y$.

We keep the complex structures on all torus fibres and on
the local finite and nodal models.  On the interpolating
annulus in B.6.16, the period matrix is smooth rather than
holomorphic in the base.  It nevertheless defines a smooth
vertical complex structure.  An invariant horizontal
distribution extends it to an almost complex structure on
the six-dimensional fibres, making the map to $B$
pseudoholomorphic wherever it is a submersion and agreeing
with the local complex structures near every critical
fibre.  Such horizontal choices form an affine space off
the critical collars, so they can be patched there and
averaged over $S^1_\delta$.  They are made on the parameter
mapping torus itself.  Thus they respect the actual closing
map, not just a chosen trivialization over an interval.

Let $E$ be the resulting complex rank-three bundle tangent
to the six-dimensional fibres of $M_s\to S^1$.  The circle
action on $E$ is complex linear, and it is free on the base.
It descends to a complex bundle $\overline E$ on $Q$.
The vector field $\delta$ and the parameter direction are
invariant trivial real lines, so the real bundles satisfy
\[
 (\overline E)_{\mathbb R}\oplus\mathbb R
                         \cong TQ\oplus\mathbb R.
 \tag{B.Q4}
\]
In particular $p_1(Q)=p_1(\overline E)$.

Write $F:Q\to B$ for the descended base map.  At a multiple
fibre of order $m_i$, divide the derivative by its known
factor of order $m_i-1$.  Its resulting target line on $Q$
is
\[
 \overline L=F^*TB\otimes
        \mathcal O_Q\!\left(-\sum_i(m_i-1)\overline S_i\right),
 \tag{B.Q5}
\]
where $\overline S_i$ is the quotient of the reduced
multiple-fibre divisor in $M_s$.  The notation
$\mathcal O_Q(\overline S_i)$ means the complex line defined
by its oriented complex normal bundle; no complex structure
on all of $Q$ is required.
The generating sphere just described lies inside the
conifold patch, whose base image is a disc.  Therefore
$F^*[\mathrm{point}]$ evaluates to zero on that sphere.
Equation (B.Q3) proves that it is zero in $H^2(Q;\mathbb Z)$.
The local root-disc model also gives, as an equality of
integral divisor classes on $Q$,
\[
 m_i[\overline S_i]=F^*[\mathrm{point}]=0.
\]
Since $H^2(Q;\mathbb Z)$ is torsion free, every
$[\overline S_i]$ is zero.  Thus
\[
                         c_1(\overline L)=0.
 \tag{B.Q6}
\]

The determinant of $\overline E$ is also trivial in
cohomology.  It suffices by (B.Q3) to test it on the new
sphere.  Its inverse image is the local vanishing $S^3$ in
\[
                         xy-zv=s\subset\mathbb C^4.
\]
The circle weights are $(1,-1,-1,1)$.  The holomorphic
residue volume form of this hypersurface is nowhere zero
and invariant: the numerator has total weight zero and
the defining equation has weight zero.  It trivializes the
equivariant determinant on this $S^3$, hence trivializes
the descended determinant on its Hopf quotient.  Therefore
\[
                         c_1(\overline E)=0.
 \tag{B.Q7}
\]
This checks the possible determinant contribution on the
new sphere; it is not inferred from a determinant computation
only on the ordinary torus locus.

Let $V$ denote the rank-two logarithmic vertical tangent
bundle.  It is the ordinary vertical tangent bundle away
from the critical curves.  Near a rank-one node with
$q=xy$ and elliptic coordinate $z$, it is the locally free
bundle with evaluation generators
$x\partial_x-y\partial_y$ and $\partial_z$.
These local extensions glue to the ordinary kernel off
the critical loci.  They descend to $\overline V$ on $Q$.
The complex vector field generated by $\delta$ is an
invariant nowhere-zero section of $V$: the circle is free,
and at either node it has a primitive component in the
surviving elliptic direction.  Hence
\[
                         c_2(\overline V)=0.
 \tag{B.Q8}
\]

After division at the multiple fibres, the derivative
has no remaining failure of surjectivity except the two
rank-one critical loci.  Near each, the complex Koszul
resolution of $(x,y)$ identifies its cokernel with the
target line restricted to the critical curve.  These local
resolutions and the actual kernel off the critical curves
give the following identity in complex topological
$K$-theory of $Q$:
\[
 [\overline E]=[\overline V]+[\overline L]
       -[\overline L|_{Z_1}]-[\overline L|_{Z_2}].
 \tag{B.Q9}
\]
Here the last two terms mean the supported $K$-classes,
extended by their complex-normal Thom classes.  Their
degree-four Chern characters are
$\operatorname{PD}_Q(Z_1)$ and
$\operatorname{PD}_Q(Z_2)$, respectively.  The identity
is consequently valid for the smooth interpolation as
well: it is exact off the critical curves, and on their
neighbourhoods it is the displayed holomorphic Koszul
calculation.  No integrability on the annulus is used.

The degree-two part of (B.Q9), together with (B.Q6)--(B.Q7),
gives $c_1(\overline V)=0$.  Its degree-four part and
(B.Q8) then give
\[
 c_2(\overline E)=\operatorname{PD}_Q(Z_1+Z_2).
\]
Equations (B.Q4) and (B.Q7) prove the first formula of
(B.Q1) integrally.  Equality first in rational cohomology
suffices here because $H^4(Q;\mathbb Z)$ is torsion free.

It remains to fix the two Euler signs.  Use the short and
long bases of B.6.16 and let $\kappa$ denote the positive
parameter direction.  At the first node the surviving
elliptic basis is $(a_2,b_2)$, its complex orientation is
that order, and
\[
 \delta=-a_2,\qquad S b_2=b_2+2a_2.
\]
With circle generator $a_2$, the critical mapping torus
has unit-line clutching
\[
 (\beta,1,z)\sim(\beta,0,e^{2\pi i\,2\beta}z).
\]
Its Chern number is $+2$ on the torus oriented by
$(b_2,\kappa)$.  For example, a unitary connection with
local form $2\pi i\,2\kappa\,d\beta$ has curvature
$-2\pi i\,2\,d\beta\wedge d\kappa$, and the convention
$c_1=iF/(2\pi)$ gives that number.  The ambient quotient
orientation with circle $\delta=-a_2$ orients $Z_1$ as
$-(b_2,\kappa)$.  Reversing both the circle and the base
orientation leaves the Euler pairing $+2$ unchanged.
At the second node the surviving basis is $(a_1,b_1)$,
$\delta=a_1$, and the return is the identity on that
lattice.  Its Euler pairing is zero.  These are the last
two assertions of (B.Q1).  Pairing its first assertion
with $x$ proves (B.Q2).
\end{proof}
